\section{Generic vanishing and nonproper arithmetic avoidance}

For the three generic-vanishing statements below, fix a commutative ring $A$, a field $K$
with a specified injective $A$-algebra map, and an $A$-module $M$.
Whenever $S$ occurs, it is a multiplicative subset of $A$ consisting of
nonzero elements.

\begin{lemma}\label{lem:kn-arithmetic-free-fibre-detection}
\lean{EnvelopeProperty.FreeFiberDetection.subsingleton_of_baseChange}
\leanok
Let $C$ be a commutative ring, $B$ a nonzero commutative $C$-algebra,
and $N$ a free $C$-module.  If $B\otimes_C N=0$, then $N=0$.
The basis cardinality need not be finite.
\end{lemma}
\begin{proof}
Scalar extension of a basis of $N$ is a basis of $B\otimes_C N$.
A basis vector over the nonzero ring $B$ is nonzero, so the hypothesis
forces the basis index set to be empty.  Coordinates in that empty basis
then identify every vector of $N$ with zero.
\end{proof}

\begin{lemma}\label{lem:kn-relative-arithmetic-torsion-base-change}
\leanok
Let $A$ be a commutative ring embedded in a field $K$, and let $M$ be
an $A$-module.  If every element of $M$ is annihilated by a nonzero
element of $A$, then $K\otimes_A M=0$.
\end{lemma}
\begin{proof}
For a pure tensor $a\otimes m$, choose $s\ne0$ with $sm=0$.
Balancing gives $s(a\otimes m)=a\otimes sm=0$.  The image of $s$ in
$K$ is nonzero by injectivity, hence invertible.  Therefore
$a\otimes m=0$.  Every tensor is a finite sum of pure tensors, so the
entire tensor product is zero.  No finiteness hypothesis on $M$ is used.
\end{proof}

\begin{theorem}\label{thm:kn-relative-arithmetic-localized-generic-vanishing}
\leanok
\uses{lem:kn-relative-arithmetic-torsion-base-change}
Let $S$ be a multiplicative subset of $A$ consisting of nonzero elements.
If $S^{-1}M=0$, then $K\otimes_A M=0$.
\end{theorem}
\begin{proof}
The equality $m/1=0$ in the localized module gives an $s\in S$ with
$sm=0$.  This $s$ is nonzero by the assumption on $S$, so the preceding
lemma applies to every $m$.
\end{proof}

\begin{theorem}\label{thm:kn-relative-arithmetic-free-fibre-generic-vanishing}
\leanok
\uses{lem:kn-arithmetic-free-fibre-detection,thm:kn-relative-arithmetic-localized-generic-vanishing}
Suppose that $S^{-1}M$ is free over $S^{-1}A$ and that $B$ is a nonzero
commutative $S^{-1}A$-algebra.  If
$B\otimes_{S^{-1}A}S^{-1}M=0$, then $K\otimes_A M=0$.
The free rank and the $A$-module $M$ need not be finite.
\end{theorem}
\begin{proof}
Choose a basis of $S^{-1}M$.  Its scalar extension is a basis of the
tensor product over $B$.  A basis vector over a nonzero ring cannot be
zero, so the basis index set must be empty.  Consequently $S^{-1}M=0$,
and the preceding theorem applies.
\end{proof}

\begin{definition}\label{def:kn-relative-arithmetic-avoiding-open}
\lean{EnvelopeProperty.FinitePresentationGenericAvoidance.avoidingOpen,EnvelopeProperty.FinitePresentationGenericAvoidance.preimage_avoidingOpen_le_compl}
\leanok
For a scheme morphism $f:X\to Y$ and a subset $Z\subseteq X$, put
$U=Y\setminus\overline{f(Z)}$.  Then $f^{-1}(U)\subseteq X\setminus Z$.
\end{definition}
\begin{proof}
The complement of the closure is open.  If $x\in Z$, then
$f(x)\in f(Z)\subseteq\overline{f(Z)}$, so $f(x)\notin U$.
This proves the inclusion without any properness hypothesis.
\end{proof}

\begin{lemma}\label{lem:kn-arithmetic-constructible-generic-closure}
\lean{EnvelopeProperty.ConstructibleGenericClosure.exists_open_deciding,EnvelopeProperty.ConstructibleGenericClosure.mem_closure_iff_of_constructible,EnvelopeProperty.ConstructibleGenericClosure.genericPoint_mem_closure_iff,EnvelopeProperty.ConstructibleGenericClosure.exists_open_avoiding}
\leanok
If $\eta$ is a generic point of the whole topological space and $Z$ is
constructible, then membership in $Z$ is constant on some open
neighbourhood of $\eta$.  In particular, $\eta\in\overline Z$ if and
only if $\eta\in Z$.  If $\eta\notin Z$, the complement of $\overline Z$
is an open neighbourhood of $\eta$ avoiding $Z$.
\end{lemma}
\begin{proof}
Induct on the construction of a constructible set from retrocompact
opens by finite unions and complements.  For an open containing $\eta$,
take that same open.  An open not containing $\eta$ is empty, because
the generic point specializes to every point; use the whole space as
the deciding neighbourhood.  For a union, intersect the deciding
neighbourhoods of the two sets.  For a complement use the same deciding
neighbourhood.  This proves the first assertion.  If $\eta\in\overline Z$,
its deciding neighbourhood meets $Z$, so its constant membership
predicate forces $\eta\in Z$.  The reverse implication follows from
$Z\subseteq\overline Z$.  Taking the complement of the closure proves
the last assertion.
\end{proof}

\begin{lemma}\label{lem:kn-relative-arithmetic-generic-avoiding-open}
\lean{EnvelopeProperty.FinitePresentationGenericAvoidance.genericPoint_mem_avoidingOpen}
\leanok
\uses{lem:kn-arithmetic-constructible-generic-closure,def:kn-relative-arithmetic-avoiding-open}
Let $Y$ be integral affine, let $f:X\to Y$ be quasi-compact and locally
of finite presentation, and let $Z$ be locally constructible.  If no
point of $Z$ lies over the generic point $\eta$ of $Y$, then $\eta\in U$.
\end{lemma}
\begin{proof}
Chevalley's theorem makes $f(Z)$ locally constructible.  Since $Y$ is
affine, hence quasi-compact, $f(Z)$ is constructible.  For a constructible
subset of an irreducible scheme, its closure contains the generic point
if and only if the subset itself does: a constructible predicate is
constant on some open neighbourhood of that point.  The hypothesis
therefore excludes $\eta$ from $\overline{f(Z)}$, proving $\eta\in U$.
\end{proof}

\begin{lemma}\label{lem:kn-relative-arithmetic-closed-constructible}
\lean{EnvelopeProperty.FinitePresentationGenericAvoidance.isLocallyConstructible_of_isClosed}
\leanok
Every closed subset of a Noetherian scheme is locally constructible.
\end{lemma}
\begin{proof}
In a Noetherian space every open is quasi-compact.  The complement of
the closed subset is therefore retrocompact open and is constructible.
Its complement is constructible and hence locally constructible.
\end{proof}

\begin{theorem}\label{thm:kn-relative-arithmetic-dense-open-avoiding-closed}
\lean{EnvelopeProperty.FinitePresentationGenericAvoidance.exists_dense_open_avoiding_closed}
\leanok
\uses{lem:kn-relative-arithmetic-generic-avoiding-open,lem:kn-relative-arithmetic-closed-constructible}
Let $Y$ be integral affine, let $X$ be Noetherian, and let $f:X\to Y$
be quasi-compact and locally of finite presentation.  If the closed
subset $Z\subseteq X$ has no point over the generic point of $Y$,
then there is a nonempty dense open $U\subseteq Y$
containing that generic point with $f^{-1}(U)\subseteq X\setminus Z$.
The morphism $f$ need not be proper.
\end{theorem}
\begin{proof}
Use the complement of $\overline{f(Z)}$.  The preceding lemmas show
that it contains the generic point and avoids the whole bad set after
inverse image.  It is nonempty because it contains that point, and a
nonempty open of the irreducible space $Y$ is dense.
\end{proof}
